\setlength{\parskip}{0.5em}
\linespread{1.5}\selectfont
\begin{adjustwidth}{0.2cm}{0.2cm}

    % Khmer Abstract
    \begin{center}
        {\khmerfont\fontsize{15pt}{25pt}\selectfont មូលន័យសង្ខេប \par}
    \end{center}
    \phantomsection
    \label{khmer-abstract}
    \vspace{0.5cm}
    \khmernormal
    \small
    បច្ចេកវិទ្យាសំយោគសំឡេងពីអត្ថបទ 
    {\englishfont\fontsize{13.5pt}{20pt}\selectfont (Text-to-Speech: TTS)} 
    សម្រាប់ភាសាខ្មែរ នៅតែជាបញ្ហាប្រឈមដ៏ធំមួយនៅក្នុងវិស័យដំណើរការភាសាធម្មជាតិ 
    {\englishfont\fontsize{13.5pt}{20pt}\selectfont (Natural Language Processing: NLP)} 
    ដោយសារតែភាសាខ្មែរត្រូវបានចាត់ទុកថាជាភាសាដែលមានធនធានទិន្នន័យតិចតួច 
    {\englishfont\fontsize{13.5pt}{20pt}\selectfont (Low-resource language)}។ 
    និក្ខេបបទនេះ បង្ហាញពីការអភិវឌ្ឍប្រព័ន្ធសំយោគសំឡេងខ្មែរដែលមានគុណភាពខ្ពស់ 
    ដោយប្រើប្រាស់ម៉ូឌែល 
    {\englishfont\fontsize{13.5pt}{20pt}\selectfont VITS (Variational Inference with adversarial learning for end-to-end Text-to-Speech)}។ 
    ការសិក្សានេះផ្ដោតសំខាន់លើការដោះស្រាយបញ្ហាសំខាន់ៗចំនួនបី៖ 
    កង្វះខាតទិន្នន័យសំឡេង និងអត្ថបទខ្នាតធំ 
    ភាពស្មុគស្មាញនៃអក្សរសាស្ត្រខ្មែរ 
    និងភាពប្រែប្រួលនៃសូរសព្ទ និងការបញ្ចេញសំឡេង។

    វិធីសាស្ត្រនៃការសិក្សានេះ ប្រើប្រាស់យុទ្ធសាស្ត្របំប្លែងអត្ថបទខ្មែរទៅជាសូរសព្ទ 
    តាមរយៈវចនានុក្រមខ្នាតធំ។ 
    ទោះបីជាការបង្ហាត់ម៉ូឌែលដោយប្រើអត្ថបទផ្ទាល់ 
    អាចផ្ដល់អត្ថប្រយោជន៍លើការចាប់យកចង្វាក់ និងរលកសំឡេងបានល្អ 
    ក្នុងករណីមានទិន្នន័យមហាសាលក៏ដោយ 
    ប៉ុន្តែសម្រាប់ការសិក្សាក្នុងស្ថានភាពដែលទិន្នន័យមានកម្រិត 
    ការប្រើប្រាស់សូរសព្ទជួយឱ្យម៉ូឌែលរៀនបានរហ័ស 
    និងមានប្រសិទ្ធភាពខ្ពស់ជាង។ 
    ប្រព័ន្ធបំប្លែងអត្ថបទទៅជាសូរសព្ទ 
    {\englishfont\fontsize{13.5pt}{20pt}\selectfont (Text-to-Phoneme: T2P)} 
    ដែលបានបង្កើតឡើង រួមមានពាក្យ និងសូរសព្ទចំនួន ៦២,១៣០ គូ 
    ដែលអាចដោះស្រាយភាពស្មុគស្មាញនៃអក្សរខ្មែរដូចជា 
    ជើងអក្សរ និមិត្តសញ្ញាផ្សេងៗ 
    និងស្រៈដែលផ្លាស់ប្តូរទីតាំង 
    ដែលជាឧបសគ្គដល់ការបំប្លែងអក្សរទៅជាសំឡេងតាមបែបបទធម្មតា។

    ទិន្នន័យដែលប្រើប្រាស់សម្រាប់ការពិសោធន៍ 
    រួមមាន ៣៣០ ប្រយោគ និងឯកសារសំឡេងចំនួន ៤០៨ ឯកសារ 
    ដែលមានរយៈពេលសរុបប្រហែល ៥៤ នាទី។ 
    ម៉ូឌែល 
    {\englishfont\fontsize{13.5pt}{20pt}\selectfont VITS} 
    រួមបញ្ចូលគ្នានូវបច្ចេកវិទ្យា 
    {\englishfont\fontsize{13.5pt}{20pt}\selectfont Variational Autoencoder (VAE)} 
    និង 
    {\englishfont\fontsize{13.5pt}{20pt}\selectfont Normalizing Flows} 
    ដើម្បីបំប្លែងលំដាប់សូរសព្ទឱ្យទៅជាសំឡេងដែលមានលក្ខណៈធម្មជាតិបំផុត។ 
    ការវាយតម្លៃត្រូវបានធ្វើឡើងតាមរយៈរង្វាស់បច្ចេកទេស 
    {\englishfont\fontsize{13.5pt}{20pt}\selectfont Word Error Rate (WER)} 
    {\englishfont\fontsize{13.5pt}{20pt}\selectfont Mel-Cepstral Distortion (MCD)} 
    និង 
    {\englishfont\fontsize{13.5pt}{20pt}\selectfont F0 Root Mean Square Error (F0 RMSE)} 
    រួមជាមួយនឹងការវាយតម្លៃតាមរយៈការស្ដាប់ដោយផ្ទាល់ 
    {\englishfont\fontsize{13.5pt}{20pt}\selectfont (Mean Opinion Score: MOS)} 
    ដើម្បីវាស់វែងភាពច្បាស់លាស់ 
    និងភាពធម្មជាតិនៃសំឡេង។

    លទ្ធផលបានបង្ហាញថា 
    ការរួមបញ្ចូលគ្នារវាងវចនានុក្រមសូរសព្ទ 
    និងម៉ូឌែល 
    {\englishfont\fontsize{13.5pt}{20pt}\selectfont VITS} 
    ផ្ដល់នូវដំណោះស្រាយដ៏មានប្រសិទ្ធភាព 
    សម្រាប់ការសំយោគសំឡេងខ្មែរ 
    ទោះបីជាក្នុងស្ថានភាពដែលមានទិន្នន័យតិចតួចក៏ដោយ។ 
    ការស្រាវជ្រាវនេះ 
    គឺជាវិភាគទានមួយក្នុងការលើកកម្ពស់អធិបតេយ្យភាពឌីជីថលសម្រាប់ភាសាខ្មែរ 
    និងជាមូលដ្ឋានគ្រឹះសម្រាប់កម្មវិធីជំនួយដល់ជនពិការភ្នែក 
    ក៏ដូចជាការផ្សព្វផ្សាយព័ត៌មានដោយស្វ័យប្រវត្តិនៅក្នុងប្រទេសកម្ពុជា។
    \vspace{2cm}
    
    \clearpage
    % English Abstract
    \begin{center}
        {\bfseries\large Abstract}
    \end{center}

    \vspace{-0.2cm}
    \setlength{\parindent}{0pt}
    \vspace{0.3cm}


    \phantomsection
    \label{abstract}
    \englishfont\fontsize{13.5pt}{17pt}\selectfont 
    ​ ​ ​ ​ ​ ​ ​​ ​ ​ ​  Text-to-Speech (TTS) technology for the Khmer language 
    remains a challenging frontier in Natural Language Processing (NLP), 
    primarily due to its classification as a low-resource language and the 
    inherent complexity of its writing system. Unlike widely supported 
    languages, Khmer suffers from a scarcity of large-scale, high-quality 
    paired text-audio datasets, which significantly limits the performance 
    of modern deep learning-based speech synthesis systems. Furthermore, 
    the Khmer abugida script introduces additional challenges, including 
    stacked consonants, overlapping diacritics, and context-dependent vowel 
    positioning, all of which complicate the direct mapping from text to 
    speech. In addition, the phonetic variability in Khmer pronunciation 
    further increases the difficulty of generating natural and intelligible 
    synthesized speech.
    
    In this work, we propose an end-to-end Khmer TTS system based on the 
    VITS (Variational Inference with adversarial learning for end-to-end 
    Text-to-Speech) architecture, specifically designed to operate under 
    low-resource constraints. To address the limitations of direct 
    text-to-speech learning, we adopt a lexicon-based Text-to-Phoneme (T2P) 
    conversion strategy, which transforms Khmer text into phoneme sequences 
    prior to model training. This approach enables more efficient learning 
    of phonetic representations and improves model convergence when data 
    availability is limited. The developed T2P system consists of 62,130 
    word-phoneme pairs and is capable of handling the structural complexity 
    of Khmer script, including stacked consonants (Cheung Akksar), multiple 
    diacritics, and positional vowels, which pose significant challenges 
    for traditional Grapheme-to-Phoneme (G2P) approaches.
    
    The experimental dataset used in this study includes 330 sentences and 
    408 high-quality audio recordings, totaling approximately 54 minutes 
    of speech data. The VITS model integrates a Variational Autoencoder 
    (VAE) and Normalizing Flows to learn a robust latent representation 
    of phoneme sequences, which are subsequently decoded into raw audio 
    waveforms using a HiFi-GAN-based generator. Model performance is 
    evaluated using objective metrics, including Word Error Rate (WER), 
    Mel-Cepstral Distortion (MCD), and F0 Root Mean Square Error (RMSE), 
    alongside subjective evaluation through Mean Opinion Score (MOS) 
    to assess naturalness and intelligibility.
    
    The experimental results demonstrate that the integration of a 
    lexicon-based phoneme representation with the VITS architecture 
    provides an effective and computationally efficient solution for 
    Khmer speech synthesis in low-resource settings. This research 
    contributes to advancing Khmer language technology and supports 
    digital sovereignty by enabling practical applications such as 
    assistive technologies for visually impaired users and automated 
    information dissemination systems in Cambodia.

    
    \end{adjustwidth}